\documentclass[twocolumn]{aastex631}
\begin{document}
\title{The reionization ionizing-photon-budget shortfall for star-forming galaxies is not robust to systematics at z\~{}5 (crisis systematic corner)}
\author{NebulaMind Lab (autonomous pipeline)}
\affiliation{NebulaMind Lab, \url{https://lab.nebulamind.net}}
\begin{abstract}
Reionization ionizing-photon-budget reconciliation at z\~{}5: star-forming galaxies require f\_esc=0.073 (+0.058/-0.032) to close the budget (Madau-Dickinson SFRD, log xi\_ion=25.5+/-0.15, clumping C=2-5, JWST-SFRD tail), versus indirect-proxy-inferred f\_esc=0.050 (+0.075/-0.030) from LzLCS O32/beta calibrations. Median delta(required-inferred)=+0.020 dex-frac (16-84\%: -0.055 to +0.083); 64\% of the systematic MC shows a shortfall. Conclusion: the budget CLOSES within the systematic, and the sign flips sign between the O32 and beta calibrations (calibration-driven, not robust). The result is bounded by the xi\_ion x clumping x proxy-calibration systematic, not by statistics. Generated autonomously from public data (jwst) via the NebulaMind Lab runner: a bounded, reproducible, descriptive study.
\end{abstract}
\section{Introduction}
Recent studies have highlighted a potential crisis in the reionization photon budget, suggesting that current observations may not account for the required number of ionizing photons to achieve reionization [Muñoz2024]. This discrepancy has led to increased scrutiny of the assumptions and calibrations used in these calculations. Previous work has emphasized the importance of accurately modeling the galaxy ionizing photon budget and its implications for reionization [Duncan2015, Davies2021].
\section{Data and method}
To address this issue, we employ a literature-anchored budget calculation that does not rely on survey catalog data. Instead, we use the Madau \& Dickinson (2014) analytic fitting function for the cosmic star formation rate density (SFRD). We adopt published values for xi\_ion and the O32/beta f\_esc proxy calibrations from LzLCS [Chisholm+22, Flury+22; Simmonds+24]. Our method focuses on reconciling systematic uncertainties in the ionizing-photon-budget through a detailed analysis of these parameters.
\section{Result}
Our reconciliation of the reionization ionizing-photon-budget at z\~{}5 reveals that star-forming galaxies require an escape fraction f\_esc=0.073 (+0.058/-0.032) to close the budget, assuming the Madau-Dickinson SFRD, log xi\_ion=25.5±0.15, and clumping factor C=2-5. In contrast, indirect-proxy-inferred f\_esc from LzLCS O32/beta calibrations yields a value of 0.050 (+0.075/-0.030). The median difference between the required and inferred escape fractions is +0.020 dex-frac (16-84\%: -0.055 to +0.083), with 64\% of systematic Monte Carlo simulations showing a shortfall.
\begin{figure}\centering\includegraphics[width=\columnwidth]{result.png}\caption{The reionization ionizing-photon-budget shortfall for star-forming galaxies is not robust to systematics at z\~{}5 (crisis systematic corner)}\end{figure}
\section{Caveats}
It is essential to acknowledge the limitations of our approach, which relies on automated, single-selection, uncalibrated measurements. The accuracy of our results depends heavily on the assumptions and calibrations used in previous studies, such as the Madau \& Dickinson (2014) SFRD and LzLCS proxy calibrations. Furthermore, our analysis does not account for potential systematic uncertainties in these underlying assumptions, which could impact the robustness of our findings. Additionally, the use of a single selection criterion may introduce biases that are not fully captured by our Monte Carlo simulations. These limitations highlight the need for further research and refined measurements to better constrain the reionization photon budget.
\section*{References}
\footnotesize
[Muoz2024] Muñoz et al. (2024). Reionization after JWST: a photon budget crisis?. bibcode 2024MNRAS.535L..37M \par [Park2022] Park et al. (2022). Calibrating excursion set reionization models to approximately conserve ionizing photons. bibcode 2022MNRAS.517..192P \par [Duncan2015] Duncan et al. (2015). Powering reionization: assessing the galaxy ionizing photon budget at z < 10. bibcode 2015MNRAS.451.2030D \par [Davies2021] Davies et al. (2021). The Predicament of Absorption-dominated Reionization: Increased Demands on Ionizing Sources. bibcode 2021ApJ...918L..35D \par [Madau2017] Madau (2017). Cosmic Reionization after Planck and before JWST: An Analytic Approach. bibcode 2017ApJ...851...50M

\end{document}
